\section{Conclusion}\label{sec:conclusion}

This paper asks what can be learned about systematic covariance when firms
are represented by distributions of risk-relevant information rather than by
point-valued characteristics.
For given latent exposure laws, quadratic transport supplies sharp upper and
lower covariance endpoints, while a nonnegative transport-excess term records
how far a realized coupling falls below the maximum-covariance arrangement.
Under the maintained transmission and return bridge, characteristic-side
distance then yields a conditional interval for the covariance ceiling.
Economically, the envelope turns distributional separation into a restriction
on feasible systematic co-movement, because, holding marginal exposure
magnitudes fixed, laws that remain far apart under every admissible
transmission face a lower covariance ceiling than nearby laws.

The empirical exercises assess whether article-embedding distance is
associated with return distance and whether the conditional restriction is
informative.
In the canonical association, one standard deviation more W$_2$ corresponds
to roughly eight fewer correlation points, and the timing-separated
later-window estimate retains the same sign.
At the same time, envelope informativeness is fragile: the zero-slack
isometric cell covers \ppnum[3]{confound-w2-l1-tau0-coverage} of observed
dyads and violates the restriction for
\ppnum[3]{confound-w2-l1-tau0-violation}, but modest slack leaves many dyads unresolved.
Together, these exercises document an association between greater
article-embedding distance and weaker return co-movement.
The rapid loss of envelope informativeness also shows how tightly the
characteristic-to-exposure map must be restricted before observed covariance
can discipline latent exposure geometry.

More broadly, distribution-valued firm information preserves within-firm
heterogeneity that point summaries omit and thereby permits economic
restrictions to be stated directly in terms of distances between probability laws.
Richer distribution-valued data could support analogous scenario analysis in
broader cross-sections, markets, or geographies, but only where the relevant
transmission and return bridges can be defended; the Nasdaq application does
not itself establish such transportability.

The application remains reduced form and associational: article embeddings
proxy latent characteristics, and the envelope is scenario analysis rather
than an identified structural estimate.
The transmission constant, exposure norms, and variance shares are maintained
rather than estimated.
The timing-separated exercise does not establish predictive validation
because the encoder has no documented training cutoff and common latent
drivers may persist across windows.
The article-embedding proxy also requires substantial firm coverage.

Within these limits, the paper shows how distribution-valued firm information
can, under explicit transmission conditions, place transparent and
economically interpretable restrictions on the feasible set of systematic covariance.
